\section{Hardness}\label{sec:complexity}

This section locates the intractability of MP-TSCFLP along the four
dimensions of an instance. We first show that shrinking two of the four sets to a single element
already forces strong NP-hardness and logarithmic inapproximability, and
then that the same covering mechanism can be lodged in the products or in
the second layer of facilities rather than in the customers. The
cardinality parameter $k$ then inherits W[2]-hardness together with a
tight ETH lower bound, and finally a purely numeric group of instances,
in which both demand dimensions are trivial, turns out to be hard as
well. The reductions grow from the second-stage
set-cover decision already visible in Section~\ref{sec:prelim}, and the computational checks behind the constructions are reproducible
from the accompanying software archive. Two conventions hold throughout
the section. In every construction with a single product the product
index is suppressed, so the data of that product are written $b_i$,
$p_j$, $q_k$, $c_{ij}$ and $d_{jk}$. We also impose a nondegeneracy
convention on the source instances, requiring the universe $U$ of each
\textsc{Set Cover} instance to be nonempty and the item total $A$ of
each \textsc{Partition} instance to be positive, for an instance
violating either requirement is a yes-instance witnessed by the empty
solution. Three recurring terms organise the section. The covering
layer, or covering mechanism, is the \textsc{Set Cover} structure a
construction embeds, and the discriminating layer is the facility side
whose cost matrix tells the members of the covered dimension apart. The
numeric layer is the hardness that survives when both demand dimensions
are trivial.

The constructions of the section share one architecture.
Theorem~\ref{thm:cover} is the canonical covering reduction,
Theorem~\ref{thm:products} and Corollary~\ref{cor:mirror} relocate the
same mechanism to the products and to the factory layer, and
Corollary~\ref{cor:coverinapprox} and Theorem~\ref{thm:w2} are
consequences of the canonical reduction rather than new constructions.
Theorem~\ref{thm:numeric} completes the fourth discrimination pair by a
different mechanism, the forced saturation of the factories, and
Table~\ref{tab:reductions} summarises the constructions side by side.

\begin{table}[H]
\small
\centering
\caption{The covering constructions of Section~\ref{sec:complexity},
one column per construction. Each column names the dimension whose
members play the \textsc{Set Cover} elements, the facility side whose
openings choose the sets, the cost block in $\{0,1\}$ that
discriminates the covered dimension, the openings forced by
feasibility, and the amplifier $Q$ that makes an uncovered element
incompatible with the budget, with $U$ the source universe, $n_U=|U|$,
$m$ the number of sets and $Q=m+1$ throughout. Theorem~\ref{thm:w2} inherits the
construction of the first column unchanged, with the cardinality bound
$k=t+1$.}
\label{tab:reductions}
\setlength{\tabcolsep}{5pt}
\begin{tabular}{@{}>{\raggedright\arraybackslash}p{2.6cm}
  >{\raggedright\arraybackslash}p{2.25cm}
  >{\raggedright\arraybackslash}p{2.25cm}
  >{\raggedright\arraybackslash}p{2.25cm}
  >{\raggedright\arraybackslash}p{2.45cm}@{}}
\toprule
 & Thm.~\ref{thm:cover} & Thm.~\ref{thm:products} &
Cor.~\ref{cor:mirror} & Thm.~\ref{thm:numeric}\\
\midrule
covered dimension & customers, $K=U$ & products, $L=U$ &
products, $L=U$ & factories, $I=U$\\[2pt]
deciding side & depots, $g\equiv 1$ & factories, $f\equiv 1$ &
depots, $g\equiv 1$ & depots, $g\equiv 1$\\[2pt]
discriminating block & second stage, $d_{ju}=0$ iff $u\in S_j$ &
first stage, $c_{i1l}=0$ iff $l\in S_i$ &
second stage, $d_{j1l}=0$ iff $l\in S_j$ &
first stage, $c_{ij}=0$ iff $i\in S_j$\\[2pt]
forced openings & the single free factory & the single free depot &
the single free factory & all $n_U$ free factories, saturated\\[2pt]
amplifier & $Q=m+1$ on each demand $q_u$ &
$Q=m+1$ on each demand $q_{1l}$ & $Q=m+1$ on each demand $q_{1l}$ &
$Q=m+1$ on each factory load\\
\bottomrule
\end{tabular}
\end{table}


\subsection{Two restrictions already suffice}\label{sec:cover}

After the preliminaries the first question is how much of the model can
be disabled before the difficulty disappears, and the answer is that
almost all of it can. One product and one factory, no factory charge,
uniform unit depot charges, and second-stage costs restricted to two
values leave a subclass that any casual reading would call trivial. The single factory feeds whichever
depots are open, so the only
decision left is which depots to install, and choosing them so that every
customer is reached through a zero-cost arc is a set-cover decision in
disguise. Figure~\ref{fig:coverage} draws the reduced instance. The
hypotheses of the next theorem are severe, and the difficulty survives all
of them at once.

% Requires: \usetikzlibrary{positioning,arrows.meta}
\begin{figure}[H]
  \centering
  \resizebox{0.8\textwidth}{!}{%
  \begin{tikzpicture}[
      >={Stealth[length=1.8mm]},
      every node/.style={font=\footnotesize},
      fac/.style={draw=black!70,line width=0.5pt,fill=black!8,rounded corners=1pt,minimum size=7mm,inner sep=1pt},
      dep/.style={draw=black!70,line width=0.5pt,circle,fill=black!8,minimum size=6.5mm,inner sep=0pt},
      cus/.style={draw=black!70,line width=0.5pt,circle,fill=black!4,minimum size=6.5mm,inner sep=0pt},
      s1/.style={draw=black!50,line width=0.35pt,->},
      cost0/.style={draw=black,line width=0.9pt,->},
      cost1/.style={draw=black!40,line width=0.35pt,->}
    ]
    % factory
    \node[fac] (I) at (0,1.2) {$I$};
    \node[font=\scriptsize,anchor=north east,align=right] at (-0.05,0.78) {$f=0$\\ huge cap.};
    % depots = sets
    \node[dep] (s1) at (3.2,3.0) {$S_1$};
    \node[dep] (s2) at (3.2,1.2) {$S_2$};
    \node[dep] (sm) at (3.2,-0.6) {$S_m$};
    \node[font=\scriptsize,anchor=south] at (s1.north) {$g=1$};
    \node[font=\scriptsize] at (3.2,0.35) {$\vdots$};
    % customers = universe
    \node[cus] (u1) at (6.6,3.0) {$u_1$};
    \node[cus] (u2) at (6.6,1.2) {$u_2$};
    \node[cus] (un) at (6.6,-0.6) {$u_{n_U}$};
    \node[font=\scriptsize] at (6.6,0.35) {$\vdots$};
    \node[font=\scriptsize,anchor=west,align=left] at (7.15,1.2) {demand\\ $Q=m+1$};
    % layer captions
    \node[font=\scriptsize\itshape] at (0,3.0) {factory};
    \node[font=\scriptsize\itshape] at (3.2,4.0) {depots $=$ sets};
    \node[font=\scriptsize\itshape] at (6.6,4.0) {customers $=$ elements};
    % stage-1 arcs cost 0
    \foreach \b in {s1,s2,sm}{\draw[s1] (I) -- (\b);}
    \node[font=\scriptsize] at (1.5,2.38) {$0$};
    % stage-2: solid cost 0 (u in S), faint cost 1 (u notin S)
    \draw[cost0] (s1) -- (u1); \node[font=\scriptsize] at (4.75,3.17) {$0$};
    \draw[cost0] (s2) -- (u2); \node[font=\scriptsize] at (4.9,1.43) {$0$};
    \draw[cost1] (s1) -- (u2); \node[font=\scriptsize] at (5.6,2.02) {$1$};
    \draw[cost1] (sm) -- (un); \node[font=\scriptsize] at (4.9,-0.36) {$1$};
  \end{tikzpicture}%
  }
  \caption{The Set Cover reduction behind Theorem~\ref{thm:cover}, with
  a single product and a single free factory of ample capacity. Each
  depot is a set $S_j$ with opening cost $1$, each customer is a
  universe element with demand $Q=m+1$, and a heavy stage-2 arc costs
  $0$ when the element lies in the set while a light arc costs $1$.
  Opening a depot is choosing its set into the cover, an uncovered
  element can be served only at cost at least $Q>B$, and the budget
  $B=t$ is therefore attainable exactly when a cover of size $t$
  exists.}
  \label{fig:coverage}
\end{figure}


\begin{theorem}\label{thm:cover}
$\mathrm{MP\text{-}TSCFLP}(B)$ is strongly NP-hard even when,
simultaneously, $|L|=1$ and $|I|=1$, the charges satisfy $f\equiv 0$ and
$g\equiv 1$, all transport costs lie in $\{0,1\}$, with in fact
$c\equiv 0$ and $d\in\{0,1\}$, and every number of the instance is bounded
by $|K|\,(|J|+1)$. In particular the problem is NP-hard already under a
unary encoding of the data.
\end{theorem}

We sketch the reduction before giving the formal proof. From an
instance of \textsc{Set Cover} with universe $U$, family
$\mathcal S=\{S_1,\dots,S_m\}$ and target $t$ we build one depot per set,
one customer per element with demand $Q=m+1$, and a single free factory
of ample capacity. The second-stage cost $d_{ju}$ is zero when
$u\in S_j$ and one otherwise, so a customer whose element is covered by
some open depot is served for free, while a customer left uncovered pays
one unit on every unit of its demand and therefore at least $Q$. Since
$Q=m+1$ exceeds both the number of depots and the target $t$, a solution can
stay within budget $B=t$ only if its open depots cover $U$, and then its
cost is exactly the number of open depots. Deciding whether budget $t$ is
attainable is thus deciding whether a cover of size $t$ exists. The
amplifier $Q$ is what makes an uncovered element incompatible with the
budget, and the per-unit accounting is what keeps
divisible service from diluting the penalty.

\begin{proof}
The plan follows the sketch, with the amplifier making the accounting
exact. We reduce from \textsc{Set Cover}, which is NP-complete
\citep{Karp1972}, under the harmless convention $1\le t\le m$. The
remaining values of $t$ are decidable in polynomial time, since an
instance with $t=0$ is a yes-instance exactly when $U=\emptyset$, while
an instance with $t>m$ is a yes-instance exactly when
$\bigcup_j S_j\supseteq U$. The nondegeneracy convention guarantees in
addition that every cover of $U$ is nonempty.

\emph{Construction.} Put $L=\{1\}$ and $I=\{1\}$ with $f_1=0$ and
capacity $b_1=n_U Q$, where $n_U=|U|$ and $Q=m+1$. Take one depot per
set, $J=\{1,\dots,m\}$ with $g_j=1$ and $p_j=n_U Q$, and one customer per
element, $K=U$ with demand $q_u=Q$. Set $c_{1j}=0$ for every $j$, set
$d_{ju}=0$ when $u\in S_j$ and $d_{ju}=1$ otherwise, and ask whether
$\OPT\le B=t$. The construction runs in $O(n_U m)$ time and its largest
number is $n_U(m+1)\le|K|(|J|+1)$.

\emph{Feasibility.} By Proposition~\ref{prop:feas} a design that opens the
factory and a depot set $Z$ is feasible if and only if $b_1\ge D_1$ and
$\sum_{j\in Z}p_j\ge D_1$, where the single product has total demand
$D_1=\sum_u q_u=n_U Q$. Since $b_1=n_U Q$ and each $p_j=n_U Q$, (F1) holds
whenever the factory is open and (F2) holds whenever $Z\ne\emptyset$.
Positive demand, guaranteed by the nondegeneracy convention, forces the
factory open, so the feasible designs are
exactly those with $Z\ne\emptyset$, and the empty design is infeasible.

\emph{Correctness.} Let $Z\subseteq J$ be the depots opened by a feasible
solution. Because $f_1=0$ and $c\equiv 0$, its cost equals
$|Z|+\sum_{j\in Z}\sum_{u\notin S_j}w_{ju}$, the transport sum ranging
over open depots because a closed depot forwards nothing
by~\eqref{eq:C4}. By constraint~\eqref{eq:C1} every unit sent to an
uncovered customer travels on an arc of cost one, so the cost is at least
$|Z|+Q\cdot V_Z$, where $V_Z$ counts the elements no open depot covers.
This lower bound holds for arbitrary divisible flows, since the second
term is charged one unit per unit of flow and an uncovered customer has
no cheaper arc available. It is attained by routing each covered
customer through a covering depot and each uncovered one through any open
depot, with the first-stage flow $x_{1j}$ set to the outflow of depot
$j$, a choice that costs nothing since $c\equiv 0$. The capacities
accommodate this routing, because $b_1=p_j=n_UQ$ equals the total demand
and so \eqref{eq:C2}--\eqref{eq:C4} hold. Now, if $C$ covers
$U$ with $|C|\le t$, then $C$ is nonempty by the nondegeneracy
convention, so opening $Z=C$ is feasible and serves every customer
for free at cost $|C|\le t$. Conversely a solution of cost at most $t$ cannot leave an
element uncovered, for that would force cost at least
$|Z|+Q\ge Q=m+1>m\ge t$, so its open depots form a cover, and
$\text{cost}\ge|Z|$ gives $|Z|\le t$. Neither direction used any
restriction beyond the budget, a fact that Theorem~\ref{thm:w2} reuses
when a cardinality bound is added.

All numbers are bounded by $n_U(m+1)$ and hence polynomial in the input
size, so the subclass remains NP-hard when its data are written in
unary, which is strong NP-hardness in the sense of
\citet{GareyJohnson1979}. A pseudo-polynomial algorithm would run in
polynomial time on this subclass, whose numbers are polynomially
bounded, and would therefore give $\Ppoly=\NP$.
\end{proof}

For coverable instances the value of the reduced instance equals the
minimum cover size, so the approximation threshold of \textsc{Set Cover} transfers verbatim.
The only care needed is that a multiplicative factor below one is
unsatisfiable once the optimum is at least one, which holds here because
a depot must open, and this is why the factor is stated as a maximum
with one.

\begin{corollary}\label{cor:coverinapprox}
For every $\varepsilon>0$, unless $\Ppoly=\NP$, no polynomial-time
algorithm approximates the optimisation version of MP-TSCFLP within a
factor $\max\{1,(1-\varepsilon)\ln|K|\}$, even with $|L|=|I|=1$,
$f\equiv 0$, $g\equiv 1$, $c\equiv 0$, $d\in\{0,1\}$ and all numbers
polynomially bounded in the instance size. The same conclusion follows under the stronger
assumption $\NP\not\subseteq\mathrm{DTIME}(n^{O(\log\log n)})$ from the
earlier bound of \citet{Feige1998}.
\end{corollary}

\begin{proof}
Enlarging the amplifier forces even approximately optimal solutions to
be covering. Rerun the construction with the
larger amplifier $Q'=n_U m+m+1$, so that $q_u=Q'$ and $b_1=p_j=n_UQ'$,
whose only effect is to widen the penalty. The cost identity of
Theorem~\ref{thm:cover}, namely that any feasible solution costs at
least $|Z|+Q'\,V_Z$ with equality for the routing described there, is
unchanged. The optimum of a coverable instance therefore equals its
minimum cover size $t^*$, attained by covering designs, since any
non-covering design costs at least $Q'$ and $Q'>m\ge t^*$ prevents it
from being optimal. Write
$\rho=\max\{1,(1-\varepsilon)\ln n_U\}$ for the claimed factor, and note
$\rho\le 1+\ln n_U$ and $t^*\le m$, so a $\rho$-approximation returns a
solution of cost at most $\rho\,t^*\le(1+\ln n_U)\,m$. Since $\ln n_U<n_U$
this is below $Q'=(n_U+1)m+1$, so the returned design is covering. The
returned solution costs at least $|Z|$ by the identity, where
$Z$ denotes its open depots, so $|Z|\le\rho\,t^*$. As
$|K|=|U|=n_U$, this yields a polynomial-time approximation of factor
$\rho$ for minimum set cover on the coverable instances. The hardness of
\citet{DinurSteurer2014} concerns exactly that minimisation problem,
whose instances are coverable by definition, since a cover is required
for the objective to exist. On instances with $(1-\varepsilon)\ln n_U<1$ the
universe has constant size at most $e^{1/(1-\varepsilon)}$, every
minimal cover has at most $n_U$ sets, and exhaustive search over the
polynomially many subfamilies of at most $n_U$ sets solves the problem
exactly in polynomial time. On all remaining instances
$\rho=(1-\varepsilon)\ln n_U$, so the combined algorithm achieves the
factor $\max\{1,(1-\varepsilon)\ln|U|\}$ on every instance of minimum
set cover, exactly on the small universes and within
$(1-\varepsilon)\ln|U|$ elsewhere. In either regime the combined
procedure would approximate minimum set cover within a factor smaller
than the threshold of \citet{DinurSteurer2014} on every instance, and
that is impossible unless $\Ppoly=\NP$. The
variant under Feige's assumption is identical.
\end{proof}

Fixing a set to a single element removes tractability entirely rather
than merely obstructing efficiency, which is the parameterized reading of the
same reduction.

\begin{corollary}\label{cor:coverpara}
$\mathrm{MP\text{-}TSCFLP}(B)$ is para-NP-hard with respect to $|L|$,
with respect to $|I|$, and with respect to the pair $(|L|,|I|)$. Unless
$\Ppoly=\NP$, none of these parameters admits an \XP\ algorithm, and a
fortiori none is \FPT.
\end{corollary}

\begin{proof}
The claims follow from the fixed slice. The hard subclass of
Theorem~\ref{thm:cover} has $|L|=|I|=1$, so a fixed
slice of each parameter is already NP-hard. A parameterized problem
NP-hard at a fixed parameter value is para-NP-hard, and an \XP\ algorithm
would decide that slice in polynomial time \citep{DowneyFellows2013}.
\end{proof}

Before the covering elements migrate to other dimensions of the model,
we record what the uniform costs of the construction do and do not imply.

\begin{remark}\label{rem:ufl}
The uniformity in Theorem~\ref{thm:cover} is deliberate. All of the
difficulty is already present with $f\equiv 0$, $g\equiv 1$ and binary
transport, so parameters counting distinct cost values, or bounding the ratio of the
positive costs, are constant on this hard subclass and help nothing on
their own. The instance is, in essence, a
non-metric uncapacitated facility location problem with a free source
stage attached, and its equivalence with \textsc{Set Cover} traces back
to \citet{Hochbaum1982}. What the theorem contributes is not the
technique but the location of the difficulty in the most restricted
subclass the model contains.
\end{remark}

\subsection{Products and customers as covering elements}\label{sec:products}

Theorem~\ref{thm:cover} kept the customers as the covering elements and
let the depots do the discriminating, but nothing forces that assignment.
If the products carry the elements and the factories discriminate through
their first-stage costs, the same amplifier argument reaches a group of
instances that Theorem~\ref{thm:cover} cannot, namely one where the
customers collapse to a single point. The hypotheses stay just as tight,
and the difficulty is now carried by the product dimension.

\begin{theorem}\label{thm:products}
$\mathrm{MP\text{-}TSCFLP}(B)$ is strongly NP-hard even when,
simultaneously, $|J|=|K|=1$, the charges satisfy $f\equiv 1$ and $g_1=0$,
first-stage costs lie in $\{0,1\}$ while second-stage costs vanish, and
every number is bounded by $|I|+1$.
The difficulty is carried by $|L|$, which remains unrestricted.
\end{theorem}

The elements have migrated from the customers to the products and
the discriminating layer from the depots to the factories, which is what
moves the hard group to $|J|=|K|=1$.

\begin{proof}
Again the source is \textsc{Set Cover} \citep{Karp1972}, with
$1\le t\le m$ and $Q=m+1$.

\emph{Construction.} Let $L=U$ and $I=\{1,\dots,m\}$ with $f_i=1$ and
$b_{il}=Q$ for every $l$. Take one depot with $g_1=0$ and $p_{1l}=Q$, one
customer with $q_{1l}=Q$, first-stage cost $c_{i1l}=0$ when $l\in S_i$
and $1$ otherwise, second-stage cost $d\equiv 0$, and budget $B=t$. The
construction takes $O(n_U m)$ time, since each cost $c_{i1l}$ is a
membership test. Full
capacity on every product makes feasibility of a design independent of
coverage, so coverage is punished only in the cost.

\emph{Correctness.} By Proposition~\ref{prop:feas}, applied per product
with $b_{il}=p_{1l}=Q$ and $D_l=q_{1l}=Q$, a design $(Y,z_1)$ is feasible
exactly when $Y\ne\emptyset$ and $z_1=1$. Products exist because
$L=U\ne\emptyset$ under the nondegeneracy convention, and for each
product (F1) reads $Q|Y|\ge Q$ and holds precisely when
$Y\ne\emptyset$, while (F2) reads $Qz_1\ge Q$ and forces $z_1=1$, an
opening that is free since $g_1=0$.
A feasible solution with open factories $Y$ and $z_1=1$ costs
$|Y|+\sum_l\sum_{i\in Y,\,l\notin S_i}x_{i1l}$, where the sum ranges over
open factories only, by~\eqref{eq:C3}. For each product the depot must receive $Q$ units
by~\eqref{eq:C1}--\eqref{eq:C2}. If no open factory contains that product
all $Q$ units travel on cost-one arcs, so the cost is at least
$|Y|+Q\,V_Y$, where $V_Y$ counts uncovered products. Equality is
attained by routing each covered product through a containing factory
and each uncovered product through any open factory, within the
per-product capacities $b_{il}=Q$, and the second stage forwards the $Q$
units to the customer at no charge. In
the forward direction a cover $C$ opens $Y=C$ and $z_1=1$ at cost
$|C|\le t$. Conversely a solution of cost at most $t<Q$ leaves no product
uncovered, since one uncovered product would force cost at least
$|Y|+Q\ge 1+Q>m\ge t$, feasibility having forced $Y\ne\emptyset$, and
then $\text{cost}\ge|Y|$ gives $|Y|\le t$, a cover. The
largest number is $Q=m+1$, so the hardness is strong
\citep{GareyJohnson1979}.
\end{proof}

The stages of the model are not interchangeable, because the balance
constraint~\eqref{eq:C2} orients the flow, so the mirror image with
depots as the discriminating layer needs its own short argument rather
than a symmetry appeal.

\begin{corollary}\label{cor:mirror}
$\mathrm{MP\text{-}TSCFLP}(B)$ is strongly NP-hard even when,
simultaneously, $|I|=|K|=1$, the charges satisfy $f_1=0$ and $g\equiv 1$,
first-stage costs vanish while second-stage costs lie in $\{0,1\}$, and
every number is bounded by $|J|+1$.
\end{corollary}

\begin{proof}
The construction of Theorem~\ref{thm:products} transfers across the two
stages, and the proof is the following list of changes. Exchange the
two facility layers, so the sets become depots $J=\{1,\dots,m\}$ with
$g_j=1$ and $p_{jl}=Q$, the single free facility becomes a factory with
$f_1=0$ and $b_{1l}=Q$, and the penalty moves onto the second stage,
$c\equiv 0$ and $d_{j1l}=0$ when $l\in S_j$ and $1$ otherwise, with
$L=U$, the demands $q_{1l}=Q$, the budget $B=t$ and the construction
time unchanged. In the correctness argument the two feasibility tests
swap parts, (F1) reading $Qy_1\ge Q$ and forcing $y_1=1$ at no charge,
(F2) reading $Q|Z|\ge Q$, that is $Z\ne\emptyset$. The cost of a
feasible solution becomes
$|Z|+\sum_l\sum_{j\in Z,\,l\notin S_j}w_{j1l}$ with the sum over open
depots by~\eqref{eq:C4}, and the lower bound $|Z|+Q\,V_Z$, with $V_Z$
the number of uncovered products, holds as before, each product's $Q$
units reaching the customer by~\eqref{eq:C1}. Equality is attained by
the routing of Theorem~\ref{thm:products} with each first-stage flow
$x_{1jl}$ set equal to $w_{j1l}$, at no cost since $c\equiv 0$, within
$b_{1l}=Q$ and satisfying \eqref{eq:C2} and~\eqref{eq:C3}. The two
directions of the equivalence then read verbatim with $Y$ replaced by
$Z$, one uncovered product forcing cost at least $1+Q>m\ge t$. Every
number of the image is again at most $Q=m+1$, so the reduction is
unary-robust and strong NP-hardness follows \citep{GareyJohnson1979}.
\end{proof}

Because these slices are strongly hard, the customer and depot parameters
already fail at the value one, and unlike the weak slice below the
difficulty does not vanish under a unary encoding.

\begin{corollary}\label{cor:parak}
$\mathrm{MP\text{-}TSCFLP}(B)$ is strongly para-NP-hard with respect to
$|K|$, to $|J|$, and to each of the pairs $(|J|,|K|)$ and $(|I|,|K|)$. No
\XP\ algorithm exists for any of these parameters unless $\Ppoly=\NP$, and
every one of these exclusions holds even under a unary encoding.
\end{corollary}

\begin{proof}
Two fixed slices supply the claims.
Theorem~\ref{thm:products} gives a strongly NP-hard slice with
$|J|=|K|=1$ and $|L|$ free, and Corollary~\ref{cor:mirror} one with
$|I|=|K|=1$. The parameterized reading is that of
Corollary~\ref{cor:coverpara}.
\end{proof}

Where no facility layer sees a demand dimension it can distinguish, the
covering mechanism becomes inactive and only the numbers remain hard. This
happens exactly with one depot and one product, and there the difficulty
drops from strong to weak.

\begin{theorem}\label{thm:partition}
$\mathrm{MP\text{-}TSCFLP}(B)$ is NP-hard even when $|J|=|L|=|K|=1$,
$g_1=0$, all transport costs vanish, and $f_i=b_i$ for every factory.
The hardness is only weak, since this group of instances is solvable in
pseudo-polynomial time $O(|I|\cdot(D+b_{\max}))$ with $D=q_{11}$ and
$b_{\max}=\max_i b_i$, and is therefore in $\Ppoly$ under a unary
encoding, hence not strongly NP-hard unless $\Ppoly=\NP$.
\end{theorem}

\begin{proof}
The reduction is from \textsc{Partition}
\citep{Karp1972}, which asks whether
integers $a_1,\dots,a_m$ with even sum $A$ admit a subset summing to
$A/2$, and the matching pseudo-polynomial program follows it.

\emph{Construction.} Take $L=K=\{1\}$ with demand $q=A/2$, one factory
per item with $f_i=b_i=a_i$, one depot with $g_1=0$ and $p_1=A/2$, all
transport costs zero, and budget $B=A/2$.

\emph{Correctness.} Every transport coefficient is zero and $g_1=0$, so
the cost of any feasible solution is $\sum_{i\in S}a_i$ with $S$ the open
factories. By Proposition~\ref{prop:feas} the design is feasible exactly
when $z_1=1$ and $\sum_{i\in S}a_i\ge A/2$. The first condition is
forced by (F2), which reads $p_1z_1=(A/2)z_1\ge A/2$ with $A>0$ by the
nondegeneracy convention, and the opening carries no charge. The second
condition is (F1). A partition $S^\ast$ opens a
feasible solution of cost exactly $A/2$, and a feasible solution of cost
at most $A/2$ has $\sum_S a_i\ge A/2$ and $\le A/2$, hence a partition. The instance
numbers are the $a_i$ themselves, written in binary, so the reduction is
not unary-robust and the hardness is only weak
\citep{GareyJohnson1979}.

\emph{The matching program.} On the general group of the statement the
depot capacity $p_1$ is arbitrary, so we first test $p_1\ge D$ in $O(1)$
time and declare the instance infeasible when the test fails. When it
holds the depot opens at no charge, the cost of opening $S$
is $\sum_{i\in S}f_i=\sum_{i\in S}b_i$, and feasibility is
$\sum_{i\in S}b_i\ge D$ by Proposition~\ref{prop:feas}, so the problem is
the knapsack-cover of minimising a subset sum subject to reaching $D$.
Let $T(i,s)\in\{0,1\}$ record whether some subset of the first $i$
factories has capacity sum exactly $s$, for $0\le s\le D+b_{\max}$. Sums
beyond $D+b_{\max}$ are irrelevant, for if a feasible $S$ has
$\sum_{i\in S}b_i\ge D+b_{\max}$ then removing any item of $S$ keeps the
sum at least $D$, so repeated removal produces a feasible subset of no
larger cost whose sum lies below $D+b_{\max}$. The recurrence
$T(i,s)=T(i-1,s)\vee T(i-1,s-b_i)$, with $T(0,0)=1$, with $T(0,s)=0$ for
$s>0$, and with the convention $T(i-1,s-b_i)=0$ when $s<b_i$, fills the
table in $O(|I|\cdot(D+b_{\max}))$ time. By induction on $i$ the
recurrence is sound and complete, every set bit corresponding to a
genuine subset sum and every subset sum in range setting its bit. The
optimum is the least $s\ge D$ with $T(|I|,s)=1$, and if no such $s$
exists the instance is infeasible, since the removal argument above
would otherwise produce a feasible sum within the range of the table.
The table
is pseudo-polynomial, so under a unary encoding this group lies in
$\Ppoly$ and the weak hardness is best possible.
\end{proof}

The same weak reduction runs with the two facility layers exchanged, which is
what closes the last quadrant of the dichotomy, the one that keeps the depots
free.

\begin{corollary}[depot mirror of the weak group]\label{cor:partitionmirror}
$\mathrm{MP\text{-}TSCFLP}(B)$ is NP-hard even when $|I|=|L|=|K|=1$, $f_1=0$, all
transport costs vanish, and $g_j=p_j$ for every depot. The hardness is only weak,
since this group is solvable in pseudo-polynomial time and hence lies in
$\Ppoly$ under a unary encoding.
\end{corollary}
\begin{proof}
Exchanging the two facility layers in the construction of
Theorem~\ref{thm:partition} yields the reduction, and the proof is the
list of changes. The items become depots with $g_j=p_j=a_j$, the single
free depot becomes a factory with $f_1=0$ and $b_1=A/2$, and the
customer, the demand $q=A/2$, the zero transport costs and the budget
$B=A/2$ are unchanged. In the correctness argument the two feasibility
tests swap parts, (F1) reading $b_1y_1=(A/2)y_1\ge A/2$ with $A>0$ by
the nondegeneracy convention, forcing the free factory opening, and
(F2) supplying the covering test $\sum_{j\in Z}a_j\ge A/2$ over the
open depots $Z$. The equivalence between solutions of cost at most
$A/2$ and partitions then reads word for word as in
Theorem~\ref{thm:partition} with $Z$ in place of $S$. The instance
numbers are the $a_j$ in binary, so the reduction is not unary-robust.
The dynamic program of Theorem~\ref{thm:partition}, run with
$g_j=p_j=a_j$ in place of $f_i=b_i$ after the $O(1)$ test $b_1\ge D$,
whose failure means infeasibility,
places the group in $\Ppoly$ under a unary encoding \citep{GareyJohnson1979}.
\end{proof}

Together the two weak slices delimit the quadrant of the model in which
the covering layer is silent, and the next observation makes that
delimitation explicit.

\begin{observation}\label{obs:quadrant}
With $|J|=|L|=1$ and $|I|,|K|$ free the covering layer is entirely
inactive. Through a single depot the factories cannot tell customers
apart, and a single product offers no dimension to discriminate either,
so the group of instances is at once weakly NP-hard, by
Theorem~\ref{thm:partition}, and pseudo-polynomially solvable. With one
product and one depot every unit of every $q_{k}$ crosses the same depot
to its customer, so the second-stage cost of any trimmed routing is the fixed constant
$\sum_k d_{1k}q_{k}$, independent of the design. Feasibility further
requires $p_1\ge\sum_k q_k$, an $O(1)$ test supplied by (F2) whose
failure means the instance is infeasible, and when the total demand is
positive the forced opening $z_1=1$ adds the constant charge $g_1$ to
every solution. What remains is a
factory-side knapsack-cover with opening charges $f_i$ and per-unit costs
$c_{i1}$ against demand $\sum_k q_k$, solvable in pseudo-polynomial time
by the factory half of the dynamic program behind
Theorem~\ref{thm:separable}, of which the subset-sum of
Theorem~\ref{thm:partition} is the special case $f_i=b_i$, $c\equiv 0$. This is the only quadrant we identify with such
a mixed profile, and it witnesses the two-layer picture directly. The
covering layer becomes active only when a facility layer sees a demand
dimension that distinguishes its members, and where none does the numeric
layer alone remains, vanishing under a unary encoding.
\end{observation}

\subsection{The parameter \texorpdfstring{$k$}{k}}\label{sec:paramk}

The genuine parameter of the study is the number of opened facilities
$k=\sum_i y_i+\sum_j z_j$ of Definition~\ref{def:Bk}, and the covering
reductions already say what to expect of it. \textsc{Set Cover}
parameterized by the cover size is the canonical W[2]-complete problem,
and the construction of Theorem~\ref{thm:cover} maps that parameter
almost onto $k$. The only subtlety is that the forced free factory
consumes one unit of cardinality without consuming budget.

\begin{theorem}\label{thm:w2}
$\mathrm{MP\text{-}TSCFLP}(B,k)$ parameterized by $k$ is W[2]-hard, even
on the subclass of Theorem~\ref{thm:cover}. The same reduction shows that
the budget-only version $\mathrm{MP\text{-}TSCFLP}(B)$ parameterized by
$B$ is W[2]-hard, and that the variant whose cardinality bound counts only depots, requiring
$\sum_j z_j\le k_z$, is W[2]-hard with parameter $k_z$, the reduction
setting $k_z=t$.
\end{theorem}

The budget parameterization is W[2]-hard here because the slice carries
unit depot charges, so the budget counts the open depots, and
parameterization by a budget value is encoding-sensitive in general.

\begin{proof}
We reduce from \textsc{Set Cover} parameterized
by $t$, which is W[2]-complete \citep{DowneyFellows2013}. The map is
that of Theorem~\ref{thm:cover},
computed in polynomial time, with $k=t+1$ a computable function of the
source parameter, so it is an \FPT\ reduction
\citep{DowneyFellows2013,Cygan2015}. For the equivalence, a cover $C$ of
size at most $t$ yields a feasible solution of cost $|C|\le t$ and
cardinality $1+|C|\le t+1$, the factory, forced open by the positive
demand, contributing the extra unit while carrying no charge. Conversely
any witness of the target side
satisfies cost at most $t$ and therefore, by the converse direction of
Theorem~\ref{thm:cover}, opens a cover of size at most $t$. The
cardinality bound only restricts the witnesses and so cannot break this
direction.

\emph{Slack.} In the image, cost at most $B$ forces
$\sum_j z_j\le\sum_j g_j z_j\le\text{cost}\le B$ because $g\equiv 1$ and
the other charges are nonnegative, while $\sum_i y_i=1$ because (F1)
with positive demand forces the single factory open. Every feasible
solution of cost at most $B$ therefore has cardinality at most $B+1$, so
the cardinality bound with $k=B+1$ is vacuous. For the budget-only version
the same instance with parameter $B=t$ is already a reduction, for its
equivalence is verbatim that of Theorem~\ref{thm:cover}, whose two
directions used only the budget. For the depot-only variant set
$k_z=t$. A cover $C$ yields a solution with $\sum_j z_j=|C|\le t$, and
the bound $k_z$ only restricts the witnesses, so the converse direction
of Theorem~\ref{thm:cover} is unaffected.
\end{proof}

The exponent of the trivial enumeration is linear in $k$, and under the
Exponential Time Hypothesis \citep{ImpagliazzoPaturiZane2001} that linear
exponent cannot be improved,
which pairs with the \XP\ membership of Observation~\ref{obs:xp} to show
that \XP\ is the right home for this parameter.

\begin{theorem}\label{thm:eth}
Assuming the ETH, no algorithm decides
$\mathrm{MP\text{-}TSCFLP}(B,k)$
in time $f(k)\cdot N^{o(k)}$, where $N$ is the instance size and $f$ is
computable, even on the subclass of Theorem~\ref{thm:cover}. Since $N$
and $n=|I|+|J|$ are polynomially related on the image, the $n^{k+O(1)}$
enumeration of Observation~\ref{obs:xp} is optimal up to the constant in
the exponent.
\end{theorem}

\begin{proof}
We compose two maps and track the parameter through the composition.
The first is the closed-neighbourhood embedding of
\textsc{Dominating Set} into \textsc{Set Cover}, whose universe is
$V(G)$ and whose sets are the closed neighbourhoods $N[v]$ for
$v\in V(G)$, so that dominating sets of size $t$ correspond exactly to
covers of size $t$. A one-vertex instance of \textsc{Dominating Set} is
decidable in constant time, so we assume $n_G\ge 2$, where $n_G$ is the
number of vertices. Composing the embedding with the map of
Theorem~\ref{thm:w2} gives
$\textsc{Dominating Set}(t)\to\mathrm{MP\text{-}TSCFLP}(B{=}t,k{=}t{+}1)$
with an image of size $N\le n_G^{c}$ for an absolute constant
$c$. The size bound follows because the
embedded \textsc{Set Cover} instance has $n_U=m=n_G$, so the image of
the map of Theorem~\ref{thm:cover} consists of an $n_G\times n_G$
second-stage cost table together with numbers bounded by $n_G(n_G+1)$.
Suppose an algorithm decided the image in
time $f(k)\cdot N^{o(k)}$. Substituting $k=t+1$ and $N\le n_G^{c}$ with $c$
constant gives $N^{o(k)}=n_G^{c\cdot o(t+1)}=n_G^{o(t)}$, since a constant
multiple of a sublinear function is sublinear and $o(t+1)=o(t)$, so the
composition would decide \textsc{Dominating Set} in time $f'(t)\cdot
n_G^{o(t)}$ for a computable $f'$, the polynomial cost of the two
reduction steps being absorbed into the product $f'(t)\cdot n_G^{o(t)}$.
This contradicts the ETH lower bound of
\citet{ChenHuangKanjXia2006}, also in \citet[Chapter~14]{Cygan2015}.
Finally, on the image the instance size $N$ and the
facility count $n=|I|+|J|$ satisfy $n\le N\le n^{c'}$ for a constant
$c'$, so an algorithm running in time $f(k)\cdot n^{o(k)}$ would run in
time $f(k)\cdot N^{o(k)}$ and is excluded by the bound just proved,
which makes the $n^{k+O(1)}$ enumeration of Observation~\ref{obs:xp}
optimal up to the constant in the exponent.
\end{proof}

The quantitative form of the exclusion requires one clarification.
Running times of the shape $f(k)\cdot N^{O(k)}$
are not, and could not be, excluded, because the enumeration of
Observation~\ref{obs:xp} achieves exactly such a bound. The excluded
class is fixed before the composition. For every computable $f$ and
every function $h$ with $h(t)/t\to 0$ no algorithm decides the target in
time $f(t)\,N^{h(t)}$, and the composition maps an algorithm with
exponent $h(k)=o(k)$ to one with exponent $h'(t)\le c\,h(t+1)+c$, which
is again $o(t)$. The pair of results thus brackets the parameter from
both sides, the exponent
linear in $k$ being attainable and, under the ETH, optimal.

\begin{remark}\label{rem:w2mem}
Membership in W[2] is left open. A witness is the set of at most $k$
opened facilities, which has the shape W[2] expects, but the acceptance
test compares binary sums and, worse, minimises an internal routing flow
by Proposition~\ref{prop:oracle}, and neither is known to translate into
bounded weft. In particular the natural verifier must evaluate a
minimum-cost flow, a computation that is P-complete in general and is
not known to be expressible by bounded-weft circuits. What can be stated without risk is that
$\mathrm{MP\text{-}TSCFLP}(B,k)$ lies in $\mathrm{W[P]}$, by guessing the
$O(k\log n)$ bits of the opening pattern and verifying in polynomial time,
so the parameter sits between W[2]-hardness and $\mathrm{W[P]}\cap\XP$.
\end{remark}

\subsection{The numeric layer}\label{sec:numeric}

The covering reductions accounted for three of the four ways a facility
layer can discriminate a demand dimension, namely depots against
customers, factories against products, and depots against products. The
remaining question is the group of instances where both demand dimensions
are trivial, $|K|=|L|=1$, and only the general transport matrix couples
the two facility layers. A computational search over small instances and candidate reduction rules,
recorded with the accompanying code, surfaced no tractable structure for this
group, and the explanation is that the group carries a fourth
discrimination pair. Tight capacities turn the factories
themselves into carriers of demand, and the first-stage matrix then
distinguishes factories that no open depot covers.

\begin{theorem}\label{thm:numeric}
$\mathrm{MP\text{-}TSCFLP}(B)$ is strongly NP-hard even when,
simultaneously, $|K|=|L|=1$, the charges satisfy $f\equiv 0$ and
$g\equiv 1$, first-stage costs lie in $\{0,1\}$ while second-stage costs
vanish, capacities are uniform on each side, and every number is bounded
by $|I|\,(|J|+1)$. Consequently this
group of instances admits no pseudo-polynomial algorithm unless
$\Ppoly=\NP$.
\end{theorem}

The discriminating pair is now the two facility layers, mediated by
tight capacities that force every feasible solution to open and
saturate all factories, and both demand dimensions stay fixed at one.

\begin{proof}
The source problem is once more \textsc{Set Cover} \citep{Karp1972},
with $1\le t\le m$ and
$Q=m+1$.

\emph{Construction.} Put $K=L=\{1\}$ with demand $D=n_U Q$, factories
$I=U$ with $f_i=0$ and $b_i=Q$, depots $J=\{1,\dots,m\}$ with $g_j=1$ and
$p_j=n_U Q$, first-stage cost $c_{ij}=0$ when $i\in S_j$ and $1$
otherwise, second-stage cost $d\equiv 0$, and budget $B=t$. As before
the construction takes $O(n_U m)$ time.

\emph{Correctness.} By Proposition~\ref{prop:feas} a design $(Y,Z)$ is
feasible exactly when (F1) $\sum_{i\in Y}b_i=Q|Y|\ge D=Qn_U$, that is
$Y=I$, and (F2) $\sum_{j\in Z}p_j\ge D$, that is $Z\ne\emptyset$, the
second equivalence using $D>0$, guaranteed by the nondegeneracy
convention. Fix a
feasible solution, write $X_j=\sum_i x_{ij}$ for the flow into depot
$j$, and write $w_j=w_{j11}$ for the flow that depot $j$ forwards to the
customer. Summing~\eqref{eq:C1} over the single customer gives $\sum_j w_j\ge D$,
summing~\eqref{eq:C2} over $j$ gives $\sum_j X_j\ge\sum_j w_j$, and
summing~\eqref{eq:C3} over $i$ gives $\sum_j X_j\le\sum_i b_i y_i=Qn_U=D$.
The chain $D\le\sum_j w_j\le\sum_j X_j\le D$ collapses to equalities, so
$\sum_j w_j=\sum_j X_j=D$ and, since each factory ships $\sum_j x_{ij}\le
b_i=Q$ while the total over the $n_U$ open factories is $Qn_U$, each
factory ships exactly $Q$. Moreover $\sum_j (X_j-w_j)=0$ with every
$X_j-w_j\ge 0$ by~\eqref{eq:C2}, so $X_j=w_j$ for all $j$. A closed depot
has $w_j=0$ by~\eqref{eq:C4}, hence $X_j=0$, so no flow enters a closed
depot and the cost sum below runs only over $j\in Z$.

The cost is then
$|Z|+\sum_i\sum_{j\in Z,\,i\notin S_j}x_{ij}$, at least $|Z|+Q\,V_Z$
with $V_Z$ the number of factories no open depot covers, since a factory
$i$ that no open depot covers ships its full load $Q$ on cost-one arcs.
Equality holds by routing each covered factory's load through a covering
depot and each uncovered factory's load through any open depot, at cost
$Q$ for the latter, then setting $w_j=X_j$ at no charge since
$d\equiv 0$, the depot capacities $p_j=n_UQ$ accommodating any inflow.
A cover $C$ opens $Y=I$, forced by (F1) and free since $f\equiv 0$,
together with
$Z=C$ at cost $|C|\le t$. Conversely the amplifier argument of
Theorem~\ref{thm:cover} again shows that a solution of cost at most $t$
leaves no factory uncovered and opens at most $t$ depots. The largest
number is $n_U(m+1)$, polynomial in the input, so the subclass stays
NP-hard when the data are written in unary \citep{GareyJohnson1979}.
\end{proof}

What survives the reduction is exactly the structure of the transport
matrix, and killing that structure restores tractability. When the matrix
is separable the two facility layers decouple and each becomes a
pseudo-polynomial knapsack-cover, so the separable matrices form a
pseudo-polynomially solvable subclass, while
Theorem~\ref{thm:numeric} makes one explicitly non-separable family
strongly hard. The classification of the territory between the two,
including structured non-separable matrices such as low-rank or Monge
matrices, remains open.

\begin{theorem}\label{thm:separable}
On the group $|K|=|L|=1$, if the first-stage matrix is separable, that is
$c_{ij}=\gamma_i+\delta_j$ for integers $\gamma,\delta$ with the arcs
nonnegative, then the problem is solvable in pseudo-polynomial time
$O\bigl(|I||J|+(|I|+|J|)\cdot D\cdot\max(b_{\max},p_{\max})\bigr)$.
Separability is decidable in $O(|I||J|)$ time.
\end{theorem}

The proof, given in \ref{app:tech}, reduces the routing cost to a
function of the two load marginals, realises any consistent pair of
loads by a northwest-corner plan, and solves each facility layer by a
dynamic program over integer loads, integrality of the split optima
coming from the vertex structure of the load polytopes. Separability is
the quadrangle condition $c_{ij}+c_{11}=c_{i1}+c_{1j}$, checked in one
pass over the matrix.

The separable subclass rounds out the numeric layer, and its
parameterized reading closes the account of the group with both demand
dimensions trivial.

\begin{corollary}\label{cor:unary}
$\mathrm{MP\text{-}TSCFLP}(B)$ is strongly para-NP-hard with respect to
the pair $(|K|,|L|)$, and unless $\Ppoly=\NP$ no \XP\ algorithm exists
for it, even under a unary encoding. Under a unary encoding the witness slices of the weak
groups become polynomial, by Theorem~\ref{thm:partition},
Observation~\ref{obs:quadrant} and Theorem~\ref{thm:separable},
while the four covering reductions of
Theorems~\ref{thm:cover}, \ref{thm:products}, \ref{thm:numeric} and
Corollary~\ref{cor:mirror} are unary-robust, so the covering layer
survives unchanged.
\end{corollary}

\begin{proof}
The claims assemble the slices already proved. The slice
$(|K|,|L|)=(1,1)$ is strongly NP-hard by
Theorem~\ref{thm:numeric}, and the parameterized reading is that of
Corollary~\ref{cor:coverpara}, now robust under a unary encoding because
the instance numbers of Theorem~\ref{thm:numeric} are polynomial. For
the remaining claims of the statement, the three positive results match
their slices as follows. The dynamic program of
Theorem~\ref{thm:partition} covers the group $|J|=|L|=|K|=1$, the
knapsack-cover of Observation~\ref{obs:quadrant} covers the quadrant
$|J|=|L|=1$, and Theorem~\ref{thm:separable} covers the separable
part of the group $|K|=|L|=1$. Each of the three runs in
pseudo-polynomial time and hence in polynomial time once the data are
written in unary. The images of the four covering reductions, in
Theorems~\ref{thm:cover}, \ref{thm:products}, \ref{thm:numeric} and
Corollary~\ref{cor:mirror}, carry numbers bounded by $|K|(|J|+1)$,
$|I|+1$, $|I|(|J|+1)$ and $|J|+1$ respectively, so all four remain valid
polynomial-time reductions when their images are encoded in unary.
\end{proof}

Taken together, the reductions of this section show that every way of
bounding a
subset $P\subseteq\{|I|,|J|,|K|,|L|\}$ that fails to retain both facility
layers, that is with $\{|I|,|J|\}\not\subseteq P$, leaves an NP-hard slice, so unless $\Ppoly=\NP$ a tractable
parameterization of this kind must retain $|I|$ and $|J|$ together. Section~\ref{sec:algorithms} takes
up exactly that surviving combination, where the enumeration of
Observation~\ref{obs:xp} sharpens into fixed-parameter tractability.
The exponent of the $n^{k+O(1)}$ enumeration itself cannot be lowered
under the ETH, by Theorem~\ref{thm:eth}.
